% problem-000440 2 活塞反应器热力学
\section*{2. 活塞反应器热力学}
如图所⽰，⽆摩擦、⽆质量的活塞1、2、3和4将反应器隔成甲、 $\mathbb { Z } _ { \setminus }$ 、丙和丁4部分，分别进⾏反应：

$
\mathrm{A} (\mathrm{g}) + \mathrm{B} (\mathrm{g}) \rightleftharpoons \mathrm{C} (\mathrm{g})
$

起始时物质的量已标在图中。某温度和100 kPa下实现平衡时，各部分的体积分别为 $V _ { \mathbb { H } }$ 、 $V _ { Z } .$ 、 $V _ { ☉ }$ 和 $V _ {  { \mathrm { T } } \circ }$ 。这时若是去掉活塞1或 $^ { 3 , }$ ，均不会引起其它活塞移动。

\textit{（原卷此处含表格/仪器试剂表，详见 PDF 或 MD 源文件。）}

\subsection*{2-1}
已知平衡常数�和反应商�，如何判断反应进⾏⽅向？

\subsection*{2-2}
求算 $\boldsymbol { { \cdot } } \boldsymbol { { V } } _ { \mathbb { H } }$ 与 $V _ { Z } ,$ 之⽐；

\subsection*{2-3}
求算�值；

\subsection*{2-4}
去掉活塞2后再次达到平衡时，活塞3向哪个⽅向发⽣了移动？试通过计算加以解释，可以假定反应的 $K ^ { \ominus }$ 等于 1。
